\documentclass[a4paper, 11pt]{article}

\usepackage[swedish]{babel}
\usepackage[utf8]{inputenc}
\usepackage[T1]{fontenc}
\usepackage{lmodern}

\usepackage{tikz}

\tikzset{
  excl/.style = {circle, draw=black, inner sep=0pt, minimum size=4pt}
}

\tikzset{
  incl/.style = {circle, draw=black, fill=black, inner sep=0pt, minimum size=4pt}
}


\usepackage{pgfplots}


\usepackage{graphicx}
\usepackage{caption}
\usepackage{subcaption}

\usepackage{enumitem}

\usepackage{multicol}

\usepackage{amsmath}
\usepackage{amsfonts}

\begin{document}

\title{
    \textbf{MM2001: Inläming 2}
}
\author{Johan Montelius}
\date{HT 2026}

\maketitle

\section*{i)}

Beräkna gränsvärdet
$$  \lim\limits_{x\to 0} \frac{e^x - e^{-x}}{sin(x)}$$
med hjälp av metoderna du lär dig på kursen 1 (det betyder särskilt:
utna l'Hospitals regel).

\begin{align*}
  \lim\limits_{x\to 0} \frac{e^x - e^{-x}}{sin(x)} &=\; 1 \\
\end{align*}



\section*{ii}

I (i) har du använt {\em räkneregler för gränsvärden.} Ange för varje
steg i din beräkning vilken regle du använder. Du kan göra det
antingen genom att ange respektive sats (samt formelnummer) i PB1
eller genom att skriva ner räkneregeln.

Det första vi gör är att förlänga uttrycket med $1/x$. Det ger oss:

$$
\lim\limits_{x\to 0}  \frac{e^x - e^{-x}}{sin(x)} =  \lim\limits_{x\to 0}  \frac{(e^x - e^{-x})/ x }{sin(x)/ x}.
$$

Detta ger oss $sin(x)/x$ i nämnaren vilket vi vet (PB 29 sid 161) har
gränsvärdet 1 då $x \rightarrow 0$. Täljaren är fortfarande ett
mysterium men vi kan skriva om uttrycket som följer:

\begin{align*}
  (e^x - e^{-x})/ x &=\; (e^x - 1 + 1 - e^{-x})/ x \\
                   &=\; (e^x - 1)/x - (e^{-x} - 1)/ x \\
\end{align*}

Vi vet (PB 29 sid 161) att gränsvärdet för $(e^x - 1)/x$ är 1 då
$x \rightarrow 0$. Gränsvärdet för den andra termen kräver lite arbete. Vi gör ett
variabelbyte, $t = -x$, vilket ger att t går mot noll då x går mot noll. Vi kan då skriva:

\begin{align*}
    \lim\limits_{x\to 0}  (e^{-x} - 1)/ x &=\;  lim{t \rightarrow 0} (e^{t} - 1)/ -t \\
                                        &=\;    \lim\limits_{x\to 0}  - (e^{t} - 1)/ t \\
                                        &=\; -1.\\.
\end{align*}

Vi har då att täljarens båda termer går mot 1 dåhar gränsvärden då $x \rightarrow 0$. 
Täljaren går mot 2 och nämnaren gick mot 1 så hela
uttrycket går mot 2.

\begin{align*}
  \lim\limits_{x\to 0}  \frac{e^x - e^{-x}}{sin(x)} &=\;  \lim\limits_{x\to 0}  \frac{(e^x - e^{-x})/ x }{sin(x)/ x} \\
                                                  &=\;  \lim\limits_{x\to 0}  \frac{(e^x - 1)/x - (e^{-x} -1)/ x}{sin(x)/ x} \\
                                                  &=\; \frac{1 - (-1)}{1} \\
  &=\; 2.
\end{align*}



\section*{iii}

Betrakta $\frac{e^{2x} - \alpha + e^{-2x}}{sin(x)}$, där $\alpha$ är ett heltal.
\begin{itemize}
\item Ange ett värde på $\alpha$ sådant att $lim_{x \rightarrow 0}\frac{e^{2x} - \alpha + e^{-2x}}{sin(x)}$ existerar. 
\item Ange ett värde på $\alpha$ sådant att $lim_{x \rightarrow 0}\frac{e^{2x} - \alpha + e^{-2x}}{sin(x)}$ inte existerar. 
\end{itemize}

Vi kan som exemple ta $\alpha = -2$ vilket ger oss möjlighet att
skriva om $e^{2x} -2 + e^{-2x}$ till $(e^x - e^{-x})^2$. 

\begin{align*}
  lim_{x \rightarrow 0}\frac{e^{2x} -2 + e^{-2x}}{sin(x)} 
      &=\; \lim_{x \rightarrow 0}\frac{(e^x - e^{-x})^2}{sin(x)} \\
      &=\; \lim_{x \rightarrow 0}\frac{e^x - e^{-x}}{sin(x)} \times (e^x - e^{-x}) \\
      &=\; \lim_{x \rightarrow 0}\frac{e^x - e^{-x}}{sin(x)} \times \lim_{x \rightarrow 0}(e^x - e^{-x}) \\
      &=\; 2 \times \lim_{x \rightarrow 0}(e^x - e^{-x}) &\; \text{från uppg. i}\\   
      &=\; 2 \times (\lim_{x \rightarrow 0}e^x - \lim_{x \rightarrow 0}   e^{-x}) \\  
      &=\; 2 \times (1 - 1) \\
      &=\; 0
\end{align*}

Om vi istället sätter $\alpha = 0$ har vi en betydligt knepigare
situation. Vi kan visserligen skriva om $e^{2x} + e^{-2x}$ till
$(e^x - e^{-x})(e^x x e^{-x})$. Men detta ger oss följande dilemma:

\begin{align*}
  lim_{x \rightarrow 0}\frac{e^{2x} + e^{-2x}}{sin(x)} 
      &=\; \lim_{x \rightarrow 0}\frac{(e^x - e^{-x})(e^x + e^{-x})}{sin(x)} \\
      &=\; \lim_{x \rightarrow 0}\frac{(e^x - e^{-x})}{sin(x)} \times \lim_{x \rightarrow 0}(e^x + e^{-x}) \\
      &=\; 2 \times \lim_{x \rightarrow 0}(e^x + e^{-x}) &\; \text{från uppg. i}.\\  
\end{align*}

Problemet här är att $e^x + e^{-x}$ inte har ett gränsvärde då
$x \rightarrow 0$. Första termen går visserligen mot 1 men andra
termen kommer gå mot $+\inf$ eller $+\inf$ beroende på om vi närmar
oss noll från höger eller vänster. Funktionen är inte kontinuerlig och
saknar gränsvärde då vi går mot noll.







\end{document}


